\section{Carrying the subsidiary's margin into the next year}
\label{sec:worked-upstream-carry}

At the first upstream year end, suppose the parent still holds the whole lot.
The subsidiary's legal profit is $81/2$: operating profit 33 plus internal
margin 10 less current tax $5/2$. Closing its legal income accounts therefore
puts $81/2$ into post-control retained profit. The parent has earned no trading
profit. These legal retained amounts cannot simply be copied into group equity.

With the inventory deferred tax asset recognised, the subsidiary's
consolidated result is $63/2$. The group close leaves $189/8$ of that result
in owners' retained profit and adds $63/8$ to acquisition NCI of $115/4$.
Opening NCI is consequently $293/8$. The account-by-account opening bridge
reconciles these closed group accounts with the closed legal books, preserving
the investment elimination, acquired basis, inventory adjustment and tax.
It does not acquire the subsidiary again. Without the inventory deferred tax
asset, the corresponding subsidiary result is 29: owners retain $87/4$ and
closing NCI is 36. The tax recognition decision has already affected the
equity that must be carried forward.

\subsection{The parent's sale releases profit attributed to the subsidiary}

In year two the parent sells the full lot outside the group for 40 cash.
Its legal books recognise cost 30 and current tax $5/2$, leaving its own
after-tax trading profit $15/2$. The group opening inventory adjustment still
reduces the parent's inventory by 10. Releasing it requires a debit to the
parent's inventory and a credit to the subsidiary's current cost of sales.
The debit removes the carried adjustment after the legal sale; the credit
locates the realised margin in the subsidiary seller's earnings. A second
charge to opening retained profit would duplicate the opening bridge.

Where the inventory deferred tax asset was recognised, its $5/2$ reversal is
also attributed to the subsidiary under the supplied premise. The subsidiary's
trading contribution is therefore $15/2$. Continue the explicit equipment
schedules: legal depreciation 14, additional acquired-basis depreciation 4,
and acquired tax-liability release 1. They contribute a net subsidiary expense
of 17, so its total second-year result is $-19/2$. NCI receives $-19/8$;
parent owners receive $3/8$ after adding the parent's own trading result.
Group profit is $-2$. The negative NCI result remains part of the allocation
\citep[paras.~B88, B94]{aasb10control2026}.

Without a recognised inventory deferred tax asset, there is no asset to reverse.
The subsidiary instead contributes margin 10 less equipment expense 17, a loss
of 7. NCI receives $-7/4$ and parent owners receive $9/4$; group profit is $1/2$.
This branch earned $5/2$ less in the first year and $5/2$ more in the second.
After full realisation, the two branches therefore reach the same closing
equity under the stipulated tax assumptions. Their yearly earnings are different.

Both branches end with cash 290, equipment 63, goodwill $135/4$, and no inventory
or inventory deferred tax asset. Assets are $1547/4$. Liabilities of $57/2$
include the original payable 20, remaining acquired deferred tax liability
$7/2$, and unpaid current taxes 5. Equity is $1433/4$, split between owners
324 and NCI $137/4$. The prior tax payable is carried; no settlement is
invented to obtain these balances.

Complete opening and closing accounts distinguish the attribution even where
aggregate profits coincide. Partial resale retains the appropriate margin and
tax; an unresolved opening or continuing tax method retains the prior result
and proposed continuation without a second-year final report.\footnote{The
upstream carry comparisons, complete expected accounts and rejection controls
are retained in
\path{data/accounting-master/acquisition-intragroup-upstream-carry-20261027/}.
Continuing control, ordinary rights, depreciation and tax schedules are supplied
conditions. Actual source and accounting admission remain unresolved.}
These acquisition examples all started with no prior interest in the subsidiary.
When an investor already holds an interest before obtaining control, the
opening investment, acquisition consideration and any remeasurement must be
distinguished before the same consolidation logic can be applied.
